\documentclass[11pt]{article}
\usepackage[margin=1in]{geometry}
\usepackage{amsmath,amssymb}
\newcommand{\eps}{\varepsilon}
\newcommand{\Fp}{\mathbb{F}_p}
\newcommand{\Q}{\mathbb{Q}}
\pagestyle{empty}
\begin{document}
\section*{How rational reconstruction works}

\paragraph{0. What is being reconstructed.}
A \emph{rational function} in $n$ variables is a ratio of two
multivariate polynomials with rational coefficients,
\[
  f(x_1,\dots,x_n) \;=\; \frac{N(x_1,\dots,x_n)}{D(x_1,\dots,x_n)},
\]
where $N$ and $D$ are each ordinary polynomials in the variables
$x_1,\dots,x_n$ whose coefficients are rational numbers (in ring
notation, $N, D \in \Q[x_1,\dots,x_n]$), with $N/D$ in lowest
terms and one denominator coefficient fixed to
$1$ to remove the overall rescaling ambiguity --- e.g.
$f = \dfrac{3x^2 + \tfrac72 xy - y}{x - 2y^2 + 5}$.
``Reconstructing'' $f$ means determining \emph{everything}: which
monomials occur in $N$ and in $D$, and every rational coefficient ---
given nothing but the ability to evaluate $f$ at chosen numerical
points.  That this is possible, and how cheaply, rests on one
observation: a single evaluation $f(t) = y$ is a \emph{linear}
condition on the unknown coefficients,
\[
  N(t) - y\, D(t) \;=\; 0 ,
\]
so roughly (\#terms of $N$) + (\#terms of $D$) probes suffice in
principle, and the whole problem is structured linear algebra plus
the number theory of step~4.

\paragraph{1. The setting in our pipeline.}
For us $f$ is a master coefficient in the variables
$C_A, C_F, \eps, x, y$: numerator degrees up to $\sim 100$ and many
thousands of terms, yet evaluable cheaply --- the trace is a black
box returning one number per point.  The strategy is therefore never
to manipulate the expression symbolically; only to sample it, and
rebuild the unique rational function consistent with the samples.

\paragraph{2. Why finite fields.}
Direct evaluation over $\Q$ blows up: numerators and denominators
grow at every step.  Instead pick a 63-bit prime $p$ and work in
$\Fp$, mapping $a/b \mapsto a\, b^{-1} \bmod p$.  Every probe is then
a fixed number of machine-word operations --- exact arithmetic with
no intermediate swell.  Probe cost depends only on the size of the
evaluation program (the trace), never on the size of the answer.

\paragraph{3. Interpolation.}
In one variable, a rational function with numerator degree $d_N$ and
denominator degree $d_D$ is fixed, up to normalization, by
$d_N + d_D + 1$ values (Thiele continued fractions, or the linear
Pad\'e system).  In many variables one interpolates recursively and
\emph{sparsely}: Zippel and Ben-Or--Tiwari methods need a probe count
scaling with the number of \emph{terms} actually present, not the
dense volume $\prod_i (d_i + 1)$.  FireFly mixes dense (one variable)
with sparse (the rest).  Hence the probe count is intrinsic to the
function --- set by its true degrees and term count --- which is why
master 107 costs $\sim$170k probes in every implementation.

\paragraph{4. Lifting back to $\Q$ (the step the name comes from).}
Interpolation delivers each coefficient as a residue
$r \bmod m$, with $m = p_1 p_2 \cdots$ assembled from several primes
by the Chinese Remainder Theorem.  \emph{Rational reconstruction}
recovers the fraction: find $a/b$ with
\[
  a \;\equiv\; b\, r \pmod m,
  \qquad |a|,\, b \;\le\; \sqrt{m/2},
\]
computed by running the extended Euclidean algorithm on $(m, r)$ and
stopping halfway (Wang's algorithm; equivalently a continued-fraction
expansion of $r/m$).  The answer is unique once $m$ exceeds twice the
square of the true numerator and denominator, so one adds primes
until the guess \emph{stabilizes}, then verifies against one more
independent prime --- the ``$2+1$ primes'' in the FireFly logs.

\paragraph{Why it is exact.}
A candidate is checked on fresh random probes.  By the
Schwartz--Zippel lemma, a wrong rational function of degree $d$
agrees with the truth at a random point of $\Fp^n$ with probability
at most $d/p \sim 10^{-14}$; every verification probe multiplies this
bound.  Probes landing on poles, or primes dividing a true
denominator, are detected and discarded.

\paragraph{Why the coefficients are rational (no $\pi$).}
Finite fields encode only fractions: $a/b \mapsto a\, b^{-1} \bmod
p$ works because $b \cdot (a/b) = a$ is an arithmetic relation that
survives modulo $p$, whereas $\pi$ satisfies no polynomial relation
with rationals at all --- it has no residue in $\Fp$, and the
Euclidean lifting can only ever return fractions.  Our coefficients
are guaranteed to be $\pi$-free because they arise from IBP
reduction: Gaussian elimination over rational functions of $d$ and
the invariants, and field operations never leave the field.  All
transcendental content ($\pi$, $\zeta$-values, polylogarithms)
lives inside the \emph{master integrals} --- which is exactly why
``hard function $=$ $\sum$ (rational coefficient) $\times$
(master)'' is the right factorization: it quarantines what
finite fields cannot represent into the objects solved by other
means.  Non-rational \emph{prefactors} that do appear
($2^{\eps}$, $\sqrt2$, $\pi^{-\eps}$-type normalizations) are
split off exactly as inert signature factors before tracing; a
constant mixed additively into a coefficient could be handled by
adjoining it as one more formal variable, at the price of an extra
variable's worth of probes.

\paragraph{The pipeline's two levers, in this language.}
Signature canonicalization reduced the \emph{number of functions} to
reconstruct (fewer columns $\Rightarrow$ fewer total probes); entry
compaction shrank the \emph{black box} (smaller trace $\Rightarrow$
cheaper probe).  The count is mathematics; the time is engineering.
\end{document}
